\documentclass[12pt, letterpaper]{article}

\usepackage[utf8]{inputenc}
\usepackage[T1]{fontenc}
\usepackage[spanish]{babel}
\usepackage{geometry}
\usepackage{amsmath, amssymb, amsthm}
\usepackage{booktabs}
\usepackage{array}
\usepackage{microtype}
\usepackage{setspace}
\usepackage{parskip}
\usepackage{titlesec}
\usepackage{fancyhdr}
\usepackage{enumitem}
\usepackage{bm}
\usepackage{mathtools}
\usepackage{tcolorbox}
\usepackage{hyperref}

\tcbuselibrary{skins, breakable}

\geometry{top=2.8cm, bottom=3.0cm, left=3.2cm, right=3.2cm, headheight=14pt}
\setstretch{1.20}

\pagestyle{fancy}
\fancyhf{}
\fancyhead[C]{\small\textsc{Econometría --- Gujarati --- Ejercicio 13.21}}
\fancyfoot[C]{\small\thepage}
\renewcommand{\headrulewidth}{0.4pt}
\renewcommand{\footrulewidth}{0pt}

\titleformat{\section}
  {\large\bfseries\scshape}{\thesection.}{0.6em}{}
  [\vspace{-4pt}\rule{\linewidth}{0.4pt}]
\titleformat{\subsection}{\normalsize\bfseries}{\thesubsection.}{0.5em}{}
\titleformat{\subsubsection}{\normalsize\itshape}{}{0em}{}

\newtcolorbox{clave}{
  colback=white, colframe=black,
  boxrule=0.5pt, arc=3pt,
  left=8pt, right=8pt, top=6pt, bottom=6pt,
  breakable
}

\newcommand{\E}{\operatorname{E}}
\newcommand{\Cov}{\operatorname{Cov}}
\newcommand{\plim}{\operatorname{plim}}

\begin{document}

\begin{center}
  {\LARGE\bfseries Error de Especificación por Variable Omitida}\\[6pt]
  {\LARGE\bfseries en la Demanda de Pollo: Pruebas RESET y LM}\\[12pt]
  {\large\itshape Detección, consecuencias de la omisión de $\ln X_{6t}$}\\
  {\large\itshape e implicaciones sobre la inclusión de precios sustitutos}\\[14pt]
  \rule{0.55\linewidth}{0.6pt}\\[8pt]
  {\normalsize Ejercicio~13.21 --- \textit{Econometría}, Gujarati \& Porter, 5.a ed.}\\[4pt]
  {\small\textsc{Análisis conceptual, teórico y algebraico}}
\end{center}

\vspace{10pt}

%------------------------------------------------------------------------------
\section{Planteamiento del problema}
%------------------------------------------------------------------------------

Se considera la \textbf{función de demanda verdadera} (modelo irrestricto):
\begin{equation}
  \ln Y_t = \beta_1 + \beta_2 \ln X_{2t} + \beta_3 \ln X_{3t}
            + \beta_6 \ln X_{6t} + u_t,
  \label{eq:verdadero}
\end{equation}
y la \textbf{función de demanda estimada} (modelo restringido, con $X_6$ omitida):
\begin{equation}
  \ln Y_t = \alpha_1 + \alpha_2 \ln X_{2t} + \alpha_3 \ln X_{3t} + v_t,
  \label{eq:omitido}
\end{equation}
donde las variables son:

\begin{center}
\renewcommand{\arraystretch}{1.3}
\small
\begin{tabular}{cl}
\toprule
\textbf{Variable} & \textbf{Definición} \\
\midrule
$Y_t$    & Consumo per cápita de pollo (libras) \\
$X_{2t}$ & Ingreso real disponible per cápita \\
$X_{3t}$ & Precio real al menudeo del pollo \\
$X_{6t}$ & Precio real compuesto de los sustitutos del pollo \\
\bottomrule
\end{tabular}
\end{center}

La especificación log-log garantiza que los coeficientes se interpretan
directamente como elasticidades. El modelo~\eqref{eq:omitido} es un
caso de \textbf{error de especificación por omisión de variable relevante}
si $\beta_6 \neq 0$ en la población.

%------------------------------------------------------------------------------
\section{Inciso a) Pruebas RESET y LM de errores de especificación}
%------------------------------------------------------------------------------

\subsection{Prueba RESET de Ramsey}

\subsubsection*{Fundamento conceptual y teórico}

La prueba RESET (\textit{Regression Equation Specification Error Test})
evalúa si existen no-linealidades omitidas o si la forma funcional del
modelo es incorrecta. Su lógica es que, si el modelo~\eqref{eq:omitido}
está bien especificado, las potencias de los valores ajustados
$\widehat{\ln Y}_t^{(2)}$ no deberían tener poder explicativo adicional
una vez que $\ln X_{2t}$ y $\ln X_{3t}$ han sido incluidas. En particular,
dado que la variable omitida $\ln X_{6t}$ está correlacionada con los
regresores incluidos, sus efectos quedarán parcialmente capturados por
las potencias de los valores ajustados.

\subsubsection*{Procedimiento algebraico}

\textbf{Paso 1.} Estimar el modelo restringido~\eqref{eq:omitido} por MCO
para obtener los valores ajustados:
\begin{equation}
  \widehat{\ln Y}_t^{(2)}
  = \hat{\alpha}_1 + \hat{\alpha}_2 \ln X_{2t} + \hat{\alpha}_3 \ln X_{3t}.
  \label{eq:yhat}
\end{equation}

\textbf{Paso 2.} Estimar la \textbf{regresión aumentada} añadiendo potencias
de los valores ajustados como regresores adicionales:
\begin{equation}
  \ln Y_t = \gamma_1 + \gamma_2 \ln X_{2t} + \gamma_3 \ln X_{3t}
            + \gamma_4 \bigl(\widehat{\ln Y}_t^{(2)}\bigr)^2
            + \gamma_5 \bigl(\widehat{\ln Y}_t^{(2)}\bigr)^3
            + \text{error}.
  \label{eq:reset_aumentado}
\end{equation}

\textbf{Paso 3.} Contrastar la hipótesis nula de correcta especificación:
\begin{equation}
  H_0: \gamma_4 = \gamma_5 = 0
  \qquad \text{vs.} \qquad
  H_1: \text{al menos uno} \neq 0.
  \label{eq:h0_reset}
\end{equation}
El estadístico $F$ de esta prueba es:
\begin{equation}
  F = \frac{(R_{\text{nueva}}^2 - R_{\text{vieja}}^2)\,/\,q}
           {(1 - R_{\text{nueva}}^2)\,/\,(n - k)},
  \label{eq:F_reset}
\end{equation}
donde $q = 2$ es el número de restricciones (términos añadidos),
$n$ es el tamaño muestral y $k$ el número de parámetros del modelo
aumentado. Bajo $H_0$, $F \sim F(q,\, n-k)$.

\textbf{Regla de decisión.} Si $F_{\text{calc}} > F_{\text{crítico}}$
al nivel $\alpha$ deseado, se rechaza $H_0$ y se concluye que el modelo
restringido~\eqref{eq:omitido} presenta un error de especificación
funcional, posiblemente debido a la omisión de $\ln X_{6t}$.

\begin{clave}
  \textbf{Interpretación de RESET en este contexto:} Si $\hat{\gamma}_4$
  o $\hat{\gamma}_5$ son significativos, los valores ajustados del
  modelo corto capturan información sistemática de $\ln X_{6t}$ que
  no fue recogida por $\ln X_{2t}$ ni $\ln X_{3t}$. Esto es evidencia
  de mala especificación en~\eqref{eq:omitido}.
\end{clave}

\subsection{Prueba del Multiplicador de Lagrange (LM)}

\subsubsection*{Fundamento conceptual y teórico}

La prueba LM de Engle (también llamada prueba de variables omitidas de
Ramsey o prueba del score) permite verificar si una o más variables
específicas deben incluirse en el modelo, a partir únicamente de los
residuales del modelo restringido. Es particularmente adecuada cuando
se conoce cuál es la variable candidata omitida (en este caso,
$\ln X_{6t}$), lo que la hace más directa que la prueba RESET para
detectar este tipo de error.

\subsubsection*{Procedimiento algebraico}

\textbf{Paso 1.} Estimar el modelo restringido~\eqref{eq:omitido} y
obtener los residuales:
\begin{equation}
  \hat{v}_t = \ln Y_t
    - \hat{\alpha}_1 - \hat{\alpha}_2 \ln X_{2t} - \hat{\alpha}_3 \ln X_{3t}.
  \label{eq:residuales}
\end{equation}

\textbf{Paso 2.} Estimar la \textbf{regresión auxiliar} de los residuales
sobre \emph{todos} los regresores del modelo verdadero~\eqref{eq:verdadero}
(los incluidos más los omitidos):
\begin{equation}
  \hat{v}_t = a_1 + a_2 \ln X_{2t} + a_3 \ln X_{3t}
              + a_4 \ln X_{6t} + e_t.
  \label{eq:reg_auxiliar}
\end{equation}

\textbf{Paso 3.} Calcular el \textbf{estadístico LM}:
\begin{equation}
  \mathrm{LM} = n \cdot R_{\text{aux}}^2,
  \label{eq:lm}
\end{equation}
donde $R_{\text{aux}}^2$ es el coeficiente de determinación de la
regresión auxiliar~\eqref{eq:reg_auxiliar} y $n$ es el tamaño muestral.

\textbf{Paso 4.} Bajo la hipótesis nula $H_0: \beta_6 = 0$ (la variable
omitida no pertenece al modelo verdadero), el estadístico LM sigue
asintóticamente una distribución ji-cuadrada:
\begin{equation}
  \mathrm{LM} \xrightarrow{d} \chi^2(m),
  \label{eq:dist_lm}
\end{equation}
donde $m = 1$ es el número de restricciones impuestas (una variable omitida).

\textbf{Regla de decisión.} Si $\mathrm{LM}_{\text{calc}} > \chi^2_{\text{crítico}}(1)$,
se rechaza $H_0$ y se concluye que $\ln X_{6t}$ pertenece al modelo
verdadero y su omisión constituye un error de especificación.

\subsubsection*{Intuición algebraica de la prueba LM}

Si el modelo restringido~\eqref{eq:omitido} fuera la especificación
correcta, los residuales $\hat{v}_t$ serían ruido blanco ortogonal a
\emph{todos} los regresores del modelo ampliado, incluyendo $\ln X_{6t}$.
Un $R_{\text{aux}}^2$ elevado en~\eqref{eq:reg_auxiliar} indica que
$\ln X_{6t}$ explica una fracción no despreciable de la variación en
$\hat{v}_t$, lo cual solo es posible si $\beta_6 \neq 0$ en la población
y la variable fue incorrectamente excluida.

\begin{clave}
  \textbf{Comparación entre RESET y LM:} La prueba RESET es más general
  (detecta cualquier mala especificación funcional) pero menos específica
  (no identifica qué variable fue omitida). La prueba LM es más directa
  cuando se conoce la variable candidata ($\ln X_{6t}$) y proporciona
  evidencia directa sobre si esa variable específica debe incluirse.
  Ambas pruebas son complementarias: si ambas rechazan $H_0$, la
  evidencia de error de especificación es robusta.
\end{clave}

%------------------------------------------------------------------------------
\section{Inciso b) ¿La no significancia de $\hat{\beta}_6$ en (1) implica
         que no hay error al ajustar (2)?}
%------------------------------------------------------------------------------

\subsection*{Respuesta: No, no lo implica}

\subsection{Argumento teórico}

La no significancia estadística de $\hat{\beta}_6$ en el modelo
irrestricto~\eqref{eq:verdadero} refleja que la muestra disponible no
proporciona evidencia suficiente para estimar con precisión $\beta_6$,
pero no implica que $\beta_6 = 0$ en la población. Existen al menos
tres razones por las que $\hat{\beta}_6$ puede resultar insignificante
aun cuando $\beta_6 \neq 0$:

\begin{enumerate}[leftmargin=1.8cm, label=\textbf{\arabic*.}]

  \item \textbf{Multicolinealidad.} Si $\ln X_{6t}$ está fuertemente
    correlacionada con $\ln X_{2t}$ o $\ln X_{3t}$, los errores estándar
    de todos los coeficientes se inflan, reduciendo las razones $t$ por
    debajo del umbral de significancia. La variable sigue siendo relevante
    pero no puede ser identificada con precisión dada la estructura de
    los datos.

  \item \textbf{Tamaño muestral reducido.} Con $n = 23$ observaciones
    anuales, la potencia de las pruebas $t$ individuales es baja.
    Un efecto de magnitud moderada ($\beta_6$ pequeño pero distinto de
    cero) puede no detectarse con tan pocos grados de libertad.

  \item \textbf{Distinción entre hipótesis nula y verdad poblacional.}
    No rechazar $H_0: \beta_6 = 0$ al nivel $\alpha = 0{,}05$ no es
    equivalente a aceptar que $\beta_6 = 0$. La ausencia de evidencia
    no es evidencia de ausencia.

\end{enumerate}

\subsection{Consecuencias algebraicas de omitir $\ln X_{6t}$ si $\beta_6 \neq 0$}

Si el verdadero modelo es~\eqref{eq:verdadero} con $\beta_6 \neq 0$ y
se estima~\eqref{eq:omitido}, el resultado del Ejercicio~13.12 se aplica
directamente. Definiendo $b_{62}$ y $b_{63}$ como los coeficientes de
la regresión auxiliar de $\ln X_{6t}$ sobre $\ln X_{2t}$ y $\ln X_{3t}$,
los estimadores MCO del modelo corto satisfacen:
\begin{align}
  \E(\hat{\alpha}_2) &= \beta_2 + \beta_6\, b_{62},
  \label{eq:sesgo_a2}\\
  \E(\hat{\alpha}_3) &= \beta_3 + \beta_6\, b_{63}.
  \label{eq:sesgo_a3}
\end{align}
Las elasticidades-ingreso ($\hat{\alpha}_2$) y precio-propio
($\hat{\alpha}_3$) son \textbf{sesgadas e inconsistentes} si
$\ln X_{6t}$ está correlacionada con los regresores incluidos, lo que
es económicamente plausible (el ingreso y los precios sustitutos
suelen moverse juntos a lo largo del ciclo económico).

\begin{clave}
  \textbf{Conclusión b):} La no significancia de $\hat{\beta}_6$ en el
  modelo completo~\eqref{eq:verdadero} no garantiza que omitir
  $\ln X_{6t}$ en~\eqref{eq:omitido} sea inocuo. Si $\beta_6 \neq 0$
  en la población pero no se detecta por multicolinealidad o muestra
  pequeña, las elasticidades estimadas en el modelo corto estarán
  sesgadas. La inferencia sobre la elasticidad-ingreso y la
  elasticidad-precio propia, que son los parámetros de mayor interés
  económico, quedará comprometida.
\end{clave}

%------------------------------------------------------------------------------
\section{Inciso c) ¿La no significancia de $\hat{\beta}_6$ implica que no
         se deben incluir precios sustitutos en la demanda?}
%------------------------------------------------------------------------------

\subsection*{Respuesta: No}

\subsection{Fundamento en la teoría microeconómica de la demanda}

La teoría de la demanda del consumidor establece que la cantidad demandada
de un bien depende de su precio, del ingreso del consumidor y de los
precios de todos los bienes relacionados (sustitutos y complementos).
Omitir los precios de los sustitutos viola el marco teórico de partida
y produce un modelo que, por construcción, no puede capturar la
respuesta de la demanda de pollo ante cambios en los precios de la carne
de res o de cerdo. Dado que la especificación log-log busca estimar
elasticidades interpretables, la exclusión de los precios sustitutos
implica que las elasticidades restantes absorben parte de un efecto
que en realidad pertenece a otra variable.

\subsection{El problema de la agregación en el índice $X_{6t}$}

La variable $X_{6t}$ es un \textbf{precio compuesto ponderado} que
agrega los precios de la carne de res y de cerdo. La no significancia
de su coeficiente puede deberse a que los efectos individuales de los
dos sustitutos se \emph{cancelan} al promediarlos en el índice:
\begin{equation}
  \ln X_{6t} = w_1 \ln P_{\text{res},t} + w_2 \ln P_{\text{cerdo},t},
  \label{eq:indice}
\end{equation}
donde $w_1$ y $w_2$ son los ponderadores del índice compuesto.
Si el precio de la res tiene una elasticidad cruzada positiva grande
y el del cerdo tiene una negativa (o viceversa), la combinación
ponderada puede producir un efecto neto cercano a cero, haciendo que
$\hat{\beta}_6 \approx 0$ aunque ambos precios individuales sean
relevantes. En este caso, la solución correcta no es eliminar los
sustitutos sino \textbf{desagregar el índice} e incluir los precios
de la res y del cerdo por separado.

\subsection{Procedimiento profesional recomendado}

Antes de concluir que los precios sustitutos no pertenecen al modelo,
el investigador debe:

\begin{enumerate}[leftmargin=1.8cm, label=\textbf{\arabic*.}]
  \item \textbf{Desagregar el índice $X_{6t}$} e incluir $\ln P_{\text{res},t}$
    y $\ln P_{\text{cerdo},t}$ por separado para detectar si los efectos
    individuales son significativos aunque el índice compuesto no lo sea.
  \item \textbf{Realizar una prueba $F$ conjunta} de la hipótesis
    $H_0: \beta_{\text{res}} = \beta_{\text{cerdo}} = 0$, en lugar de
    confiar únicamente en pruebas $t$ individuales que pueden ser
    débiles bajo multicolinealidad.
  \item \textbf{Consultar la teoría económica}: la demanda de pollo
    debe depender de los precios de los sustitutos cercanos; la evidencia
    contraria requiere una justificación sustantiva, no solo estadística.
\end{enumerate}

\begin{clave}
  \textbf{Conclusión c):} La no significancia de $\hat{\beta}_6$
  no implica que los precios sustitutos deban excluirse del modelo de
  demanda. Puede reflejar un problema de agregación en el índice compuesto,
  multicolinealidad o baja potencia de la prueba individual. La solución
  correcta es desagregar el índice, aplicar una prueba de restricciones
  conjuntas y respetar el marco teórico microeconómico que predice la
  relevancia de los precios sustitutos en cualquier función de demanda
  bien especificada.
\end{clave}

\end{document}
